% !TeX root = ../papers/nc-p3.tex
\section{有理正規三次上の六次二重曲線の全排除}
\label{sec:double-twisted-cubic-exclusion}
\begin{theorem}\label{thm:double-twisted-cubic-exclusion}
一次元の点ファイバーグラフを持つ標準有限Frobenius窓において、純部分が
空間有理正規三次の上の重複度2の曲線となることはない。
\end{theorem}
\begin{proof}[二重曲線の抽象型]
定理\ref{thm:pure-curve-defect}より点グラフ $X$ 自体が純で $p_a=3$ である。
被約台を $C=\mathbb P^1$ とし、そのイデアルを $N$ とする。
$N^2$ は一般点で零、純性から全体で零である。$N$ は滑らかな $C$ 上で階数1無捩れ、
Euler標数から $N=\mathcal O_C(-4)$ である。このような二重曲線（リボン）の抽象型は二つしかない。
二つのアフィン座標直線では各パラメータを持ち上げて構造層の商写像を環として分裂できる。
重なりでの分裂の差は $N$ に値を持つ導分で、変更を除いた全体は
$H^1(T_C\otimes N)=H^1(\mathcal O(-2))=k$、$N$ の倍率を許すと零と非零の二型である。

分裂型では引戻し $J=f^*\mathcal O(1)$ を持つ有限平坦二重被覆 $f:X\to C$ がある。
その相対双対は $\operatorname{Hom}(\mathcal O\oplus\mathcal O(-4),\mathcal O(-2))$ なので
$\omega_X=J^2$ である。各 $L_i$ の制限は $\mathcal O_C(3)$ で、$N L_i|_C=\mathcal O(-1)$ の
コホモロジー消滅から全四切断になる。$L_i/J^2$ が非零切断を持てば、冪零部分の次数は $-3$ なので
一般点で非零因子であり、長さ2のCartier零因子を持つ。その上で $L_i$ の三切断が零になり埋込みに反する。
ゆえに $f_*L_i=\mathcal O(1)^2$。$P_i=L_i/J$ の二切断は各長さ2ファイバーを張る。
これで定理\ref{thm:smooth-sextic-exclusion}の超楕円分裂型の証明を適用できる。
同じ $P_i=P_{i+1}$ ならば二鉛筆の積の滑らかな二次曲面の対角が入る。
実際 $X$ を含む二次曲面はこの場合一意で、独立な二式は次数6を切れず、
共通平面因子を持てば非平面な台が直線に入ってしまう。
$P_i\ne P_{i+1}$ の場合の必要な次数零切断の消滅も、冪零部分の負次数と
$C$ 上の定数を用いれば従う。他は有限平坦な各長さ2の環の乗法だけを使うので同じ三次元数上限18を得る。

非分裂型は平面内の二重円錐曲線 $(xz-y^2)^2=0$ である。
実際 $x=1$ の座標を $t=y/x,e=z/x-t^2$、$z=1$ の座標を $u=y/z,e'=x/z-u^2$ とすると
\[
                 e^2=0,\qquad u=t^{-1}-e t^{-3},\quad e'=e t^{-4}.
\]
この分裂の差は $e\partial_t$ の自明化で $t^{-1}$ に対応し、上の $H^1(\mathcal O(-2))$ の生成元である。
随伴によりこの平面埋込みが完全標準系であり、全自己同型は標準平面の二次曲線安定群 $\mathrm{PGL}_2$ である。

\smallskip
\noindent\textit{四移動と全切断系。}
$1+N$ の単元完全列より $\operatorname{Pic}^0 X=H^1(N)$ は三次元ベクトル群である。
$H^1(\mathcal O(-4))$ への $\mathrm{PGL}_2$ 作用は二項二次形式の三次元表現である。
実際、Serre双対で $H^0(\mathcal O(2))$ の双対、$N$ の線型化は
標準平面二次曲線の余法束によるものなので、半単純群に指標の違いもない。
被約次数3の不変な参照束として、$\omega_X^3$ の一意な平方根を取れる
（被約上の平方根を持ち上げてベクトル群で2を割る）。これとの差を $l_i\in H^1(N)$ と書く。
偏極三階差の消滅から $l_i$ は $i$ の二次以下の多項式であり、$l_{i+4}=\sigma^*l_i$。
$\sigma$ が半単純なら各固有成分でこの等式を見て多項式は定数である。
二次曲面に入らないことは下記の座標でも示すので、この場合全隣接関係が交換子となり点グラフは全空間になる。
恒等もここに含める。
非自明ユニポテントなら同じ一径数加法群内の四乗根を $\tau$ として
$l_i=(\tau^*)^i l_0$ となる。実際右辺も三次元表現の冪零指数から二次以下の多項式で、
$0,4,8$ で一致する。従って全四切断系は一つの系と連続する $\tau$ の移動で与えられる。

各束 $L$ の遷移函数は二つのアフィン座標で
$t^3(1+e(a/t+(b+3/2)/t^2+c/t^3))$ と書ける。
三係数は $H^1(\mathcal O(-4))$ のLaurent基底であり、$3/2$ は基準の $\omega_X^{3/2}$ による。
$t^j$ の一意な切断持上げは $t^j+e g_j(t)$ で
\[
 (g_0,g_1,g_2,g_3)=(0,a,at+b+1/2,at^2+(b+3/2)t+c).
\]
これは $u^{3-j}$ を先の遷移式で書き、$e$ の係数の非負べきだけを残せば得られる。
三次曲線の三つのHankel二次式に代入した $e$ 係数は、斉次化して
$sU,tU+sV,tV$、$U=(b+1/2)s-at$、$V=cs+(1/2-b)t$ である。
従って埋込みであるためには
\[
                          1/4-b^2+ac\ne0
\]
が必要である。さもなければ共通零点の法方向が三次曲線の接方向に入る。
この不等式で三つの二次係数は独立で、二重曲線を含む二次曲面はない
（三係数の行列式が同じ $1/4-b^2+ac$ である）。

二つの隣接偏極が異なる場合にも、六関係は全切断乗法核である。
被約台での七次元の二文字像に加え、冪零部分の $H^0(\mathcal O(2))$ も全て得られる。
二つの束の上記表示の係数差を $A,B,C'$ とする。交換子の $e$ 係数には
\[
 A,\quad At+B,\quad At^2+Bt+C',\quad Bt,\quad Bt^2+C't,\quad C't^2
\]
があり、三係数の少なくとも一つが非零ならこれらは二次以下の多項式を張る。
従って像次元10、核次元6である。同じ偏極の場合にも先の二次係数で同じ結論となる。
従って $\tau$ で系を運べば全辺は同一の通常関係であり、$j_i=j_0\tau^i$ と選べる。
$\tau^4=\sigma$ なので、この通常表示での位置4から位置0へのFrobenius同定はスカラーである。
実際その射影変換は両側の $j_4=j_0\sigma$ の同じ座標を同一視し、非退化な被約三次を点ごとに固定する。
従って高さ4のテンソルは通常の四位置巡回の固有ベクトルである。固有値を先に指定する必要はない。

\smallskip
\noindent\textit{巡回テンソルの全係数排除。}
標準平面の座標を取り $\tau(t,e)=(t+1,e)$ とできる。四次テンソルを
四つの変数 $t_j$ ($j$ は法4) で各次数3以下の多項式 $w$ と見る。
全辺で $t_{j+1}=t_j+1$ の三次曲線のグラフ上に消えるため
\[
                  w=\prod_j(t_{j+1}-t_j-1)\, f(t_0,t_1,t_2,t_3).
\]
$0\ne f$ は各次数1以下であり、前の積が巡回不変なので $f$ も巡回固有である。
最初の辺がリボン全体を消すための一次の必要条件は
\[
       [(D_0+D_1)w]_{t_0=t,\ t_1=t+1}=0,\qquad D(t^k)=g_k(t).
\]
これは他の二位置を被約台に制限しても全四文字を張るので係数恒等式となる。

使用した消去を再現できる形で示す。$\{0,1\}^4$ の巡回軌道を辞書順に取り、
固有値の候補 $z=1,-1,\mathrm i,-\mathrm i$ に対し長さ $h$ が $z^h=1$ を満たす軌道のみ採用する。
その単項式を順に $z^j$ 倍して足した列を使う。各列に四つの括弧を掛け、位置0,1の指数
$k$ を $m=1,2,3$ だけ下げる係数を順に $2a,2b+2k-3,2c$ とする（$k<m$ は零）。
$(t+1)^d$ を展開して $t,t_2,t_3$ の係数を並べれば上式の二倍である。
軌道係数の最初の非零項を1にした各射影チャートで、この行列の消滅イデアルを
$\mathbb Q(\mathrm i)$ 上で計算した。$z=-1,\mathrm i$ の全 $4,3$ チャート、および
$z=1$ の最初以外の五チャートは単位イデアルである。共役で $-\mathrm i$ も同じである。
$z=1$ の最初のチャートのイデアルは $a^2,(b+1/2)^2$ を含む。
この計算は\path{computations/verify_double_twisted_cubic.py}が全行と各Gröbner基底を構成し、
二多項式の剰余と単位イデアルを検証する。最後の場合も $a=0,b=-1/2$ となり
埋込みの不等式に反する。これで全ての二重構造を排除した。
\end{proof}
